\documentclass[11pt]{article}
\usepackage[margin=1in]{geometry}
\usepackage{booktabs}
\usepackage{array}
\usepackage{tabularx}
\usepackage{longtable}
\usepackage{titlesec}
\usepackage{parskip}
\usepackage{microtype}
\usepackage{hyperref}
\usepackage{enumitem}

\hypersetup{colorlinks=true, linkcolor=black, urlcolor=blue}

\titleformat{\section}{\large\bfseries}{}{0em}{}[\titlerule]
\titleformat{\subsection}{\normalsize\bfseries}{}{0em}{}

\title{\textbf{Burst Splitter Internal Logic}}
\date{}
\author{}

\begin{document}
\maketitle
\thispagestyle{empty}

\section{Overview}

The Burst Splitter is a two-stage pipelined block shared across all \texttt{N\_CLIENTS} AXI clients. It sits between the AW/AR FIFO and the Packet Formation Unit. Its function is to convert validated AXI burst requests into normalized per-beat packets, splitting bursts that cross DRAM row boundaries into two sub-transactions.

\medskip
\noindent\textbf{Input:} AW/AR FIFO output --- \texttt{ADDR}, \texttt{AXID}, \texttt{AXLEN}, \texttt{AXSIZE}, \texttt{AXBURST}

\noindent\textbf{Output:} Per-beat context (\texttt{addr}, \texttt{tag}, \texttt{pkt\_type}, \texttt{beat\_cnt}) $\rightarrow$ Packet Formation Unit via \texttt{pkt\_valid}

\section{Stage 1 --- Decode \& Compute (1 cycle, registered output)}

Stage 1 runs combinatorially and captures all results into the S1$\rightarrow$S2 pipeline register on the next clock edge. Three sub-blocks operate in parallel:

\subsection{Burst Type Decoder}
\begin{itemize}[noitemsep]
  \item Input: \texttt{AXBURST[1:0]} from FIFO
  \item Output: \texttt{is\_wrap} / \texttt{is\_incr} flag $\rightarrow$ pipeline register
  \item FIXED burst type is not supported (design-time assertion)
\end{itemize}

\subsection{Row Boundary Checker}
\begin{itemize}[noitemsep]
  \item Inputs: \texttt{ADDR}, \texttt{AXLEN}, \texttt{AXSIZE} from FIFO
  \item Computes: \texttt{addr\_start + burst\_bytes} vs \texttt{ROW\_MASK}
  \item Outputs: \texttt{split\_needed} flag, \texttt{row\_boundary} address $\rightarrow$ Split Point Calc and Beat Count Calc
\end{itemize}

\subsection{WRAP Address Generator}
\begin{itemize}[noitemsep]
  \item Inputs: \texttt{ADDR}, \texttt{AXLEN}, \texttt{AXSIZE} from FIFO
  \item Active only when \texttt{is\_wrap = 1}
  \item Outputs: \texttt{wrap\_boundary}, \texttt{beat\_addrs[0..N]} $\rightarrow$ pipeline register
  \item Maximum 16 entries (AXI4 WRAP burst max \texttt{AXLEN} = 15)
\end{itemize}

\subsection{Split Point Calc}
\begin{itemize}[noitemsep]
  \item Inputs: \texttt{row\_boundary} from Row Boundary Checker; \texttt{ADDR} from FIFO
  \item Computes: \texttt{addr\_b = row\_boundary}, \texttt{beat\_cnt\_b = total\_beats $-$ beat\_cnt\_a}
  \item Output: \texttt{addr\_b}, \texttt{beat\_cnt\_b} $\rightarrow$ pipeline register
\end{itemize}

\subsection{Beat Count Calc}
\begin{itemize}[noitemsep]
  \item Inputs: \texttt{row\_boundary} from Row Boundary Checker; \texttt{AXLEN} from FIFO
  \item Computes: \texttt{beat\_cnt\_a = beats\_to\_row\_end}, \texttt{total\_beats = AXLEN + 1}
  \item Output: \texttt{beat\_cnt\_a}, \texttt{total\_beats} $\rightarrow$ pipeline register
\end{itemize}

\subsection{S1$\rightarrow$S2 Pipeline Register (clocked)}

\begin{center}
\texttt{\{ addr\_a, beat\_cnt\_a, addr\_b, beat\_cnt\_b, split\_needed, is\_wrap, beat\_addrs[], axid, pkt\_type, tag \}}
\end{center}

\section{Stage 2 --- Emit \& Track}

Stage 2 is multi-cycle when a split is needed. It stalls Stage 1 (via \texttt{s1\_busy}) while emitting sub-transaction B.

\subsection{Emit FSM}

\begin{itemize}[noitemsep]
  \item \texttt{IDLE $\rightarrow$ EMIT\_A}
  \item \texttt{EMIT\_A $\rightarrow$ IDLE} \hfill (split\_needed = 0)
  \item \texttt{EMIT\_A $\rightarrow$ EMIT\_B} \hfill (split\_needed = 1)
  \item \texttt{EMIT\_B $\rightarrow$ IDLE}
  \item Hold current state if \texttt{pkt\_ready = 0}
  \item Asserts \texttt{s1\_busy} during \texttt{EMIT\_B} to stall Stage 1
\end{itemize}

\subsection{Sub-transaction A Mux}
\begin{itemize}[noitemsep]
  \item Selects: \texttt{addr\_a}, \texttt{beat\_cnt\_a}, \texttt{tag}, \texttt{pkt\_type} from pipeline register
  \item Forwards to Pkt Assembler on cycle N
\end{itemize}

\subsection{Sub-transaction B Mux}
\begin{itemize}[noitemsep]
  \item Selects: \texttt{addr\_b}, \texttt{beat\_cnt\_b}, same \texttt{tag}, \texttt{pkt\_type} from pipeline register
  \item Forwards to Pkt Assembler on cycle N+1
\end{itemize}

\subsection{Beat Tracker}
\begin{itemize}[noitemsep]
  \item Maintains \texttt{beat\_cnt\_reg}; decrements on \texttt{pkt\_accept}
  \item Asserts \texttt{last\_beat} flag when count reaches zero
  \item Output: \texttt{last\_beat} $\rightarrow$ Pkt Assembler
\end{itemize}

\subsection{WRAP Beat Address Table}
\begin{itemize}[noitemsep]
  \item Up to 16-entry register file, written once on S1$\rightarrow$S2 pipeline register transfer
  \item Input: \texttt{beat\_addrs[]} from S1$\rightarrow$S2 pipeline register
  \item Read sequentially by \texttt{beat\_idx} counter
  \item Output: current beat address $\rightarrow$ Pkt Assembler
\end{itemize}

\subsection{Pkt Assembler}
\begin{itemize}[noitemsep]
  \item Inputs: \texttt{addr}, \texttt{tag}, \texttt{pkt\_type}, \texttt{beat\_cnt} from Sub-txn A/B mux; \texttt{last\_beat} from Beat Tracker; beat address from WRAP beat addr table
  \item Assembles normalized read (42-bit) or write (618-bit) packet
  \item Asserts \texttt{pkt\_valid} $\rightarrow$ Packet Formation Unit
\end{itemize}

\section{Backpressure \& Control Signals}

\begin{center}
\begin{tabularx}{\textwidth}{>{\ttfamily}l >{\ttfamily}l X}
\toprule
\textbf{Signal} & \textbf{Direction} & \textbf{Description} \\
\midrule
s1\_busy   & Emit FSM $\rightarrow$ AW/AR FIFO     & Stalls Stage 1 from consuming next FIFO entry during EMIT\_B \\
pkt\_ready & Packet Formation Unit $\rightarrow$ Emit FSM & Back-pressure from downstream; FSM holds state when low \\
pkt\_valid & Pkt Assembler $\rightarrow$ Packet Formation Unit & Indicates normalized packet ready for consumption \\
\bottomrule
\end{tabularx}
\end{center}

\section{Three Paths Through the Splitter}

\begin{description}[noitemsep]
  \item[Path A --- No split (INCR, no row crossing)] \hfill \\
    \texttt{split\_needed = 0}. Stage 2 emits one packet (\texttt{EMIT\_A $\rightarrow$ IDLE} in 1 cycle). \texttt{s1\_busy} never asserted. Throughput: 1 transaction / cycle.

  \item[Path B --- Row-boundary split (INCR, crosses boundary)] \hfill \\
    \texttt{split\_needed = 1}. Stage 1 computes \texttt{addr\_b}, \texttt{beat\_cnt\_b}. Stage 2: \texttt{EMIT\_A} cycle N, asserts \texttt{s1\_busy}, \texttt{EMIT\_B} cycle N+1. Throughput: 1 transaction / 2 cycles.

  \item[Path C --- WRAP burst] \hfill \\
    WRAP address generator fills \texttt{beat\_addrs[]}. Stage 2 iterates via WRAP beat addr table indexed by \texttt{beat\_idx}. May also split at row boundary.
\end{description}

\section{Internal State}

\begin{center}
\begin{tabular}{>{\ttfamily}l l}
\toprule
\textbf{Field} & \textbf{Description} \\
\midrule
base\_addr   & Base address of current burst \\
beat\_idx    & Current beat index within burst \\
beats\_left  & Remaining beats to emit \\
burst\_type  & INCR or WRAP \\
burst\_len   & Total beats (AXLEN + 1) \\
txn\_tag     & Allocated transaction tag \texttt{[AXID[5:0] + seqnum[1:0]]} \\
\bottomrule
\end{tabular}
\end{center}

\section{Timing Notes}

\begin{itemize}[noitemsep]
  \item Stage 1 combinatorial depth must close timing in 1 cycle. WRAP address generator is the critical path --- limited to AXLEN $\leq$ 15 for WRAP bursts (AXI4 constraint).
  \item S1$\rightarrow$S2 register clocked at \texttt{posedge clk}.
  \item \texttt{pkt\_ready = 0} holds the Emit FSM in its current state; \texttt{s1\_busy} remains asserted until \texttt{IDLE} is reached.
  \item Tag Allocator must have a valid tag before Stage 1 can proceed --- Tag Allocator stall folds into \texttt{s1\_busy}.
  \item Error path (\texttt{pkt\_type = Err\_write / Err\_read}): set in Stage 1, carried through pipeline register unchanged. Both sub-transactions on a split carry the same \texttt{pkt\_type}.
\end{itemize}

\end{document}
